% block: hermitian-linear-systems-heading page:85 type:section
\subsection{\BT{Hermitian Linear Systems}{Hermitian 线性方程组}}\label{sec:hermitian-linear-systems}
\BT{In this section, we give derivations of the Conjugate Gradient (CG) method, the Conjugate Residual (CR) method, and the Minimal Residual (MINRES) method. These methods are used for the case where matrix $A$ is Hermitian (or real symmetric). The CG method is usually used for the case where $A$ is a Hermitian positive definite matrix, and the CR method and the MINRES method are used for a Hermitian (indefinite) matrix. In exact precision arithmetic, the CR method and the MINRES method produce the same approximate solutions. We will see that these methods find the best approximate solutions at each iteration step. To be specific, the CG method for a Hermitian positive definite matrix finds the best approximate solution such that a weighted norm of the error is minimized. On the other hand, the CR method and the MINRES method find the best approximate solution such that the 2-norm of the residual is minimized.}{本节推导共轭梯度 (CG) 方法, 共轭残差 (CR) 方法和最小残差 (MINRES) 方法. 这些方法用于矩阵 $A$ 为 Hermitian (或实对称) 矩阵的情形. CG 方法通常用于 $A$ 为 Hermitian 正定矩阵的情形, CR 方法和 MINRES 方法则用于 Hermitian (不定) 矩阵. 在精确算术下, CR 方法和 MINRES 方法产生相同的近似解. 我们将看到, 这些方法在每次迭代都找到最佳近似解. 具体而言, 对 Hermitian 正定矩阵, CG 方法找到使误差的某个加权范数达到极小值的最佳近似解. 另一方面, CR 方法和 MINRES 方法找到使残差的 2-范数达到极小值的最佳近似解.}
% block: cg-heading pages:85,86 type:subsection
\subsubsection{\BT{The Conjugate Gradient (CG) Method}{共轭梯度 (CG) 方法}}\label{subsec:conjugate-gradient}
\BT{The CG method \BookCite{95} was proposed by Hestenes (1906--1991) and Stiefel (1909--1978) in 1952, and is the best-known Krylov subspace method. In exact precision arithmetic, as well as other Krylov subspace methods the CG method produces the exact solution within finite iteration steps. This means that the CG method has two features: a direct method as in Sect.\ \ref{sec:direct-methods} and an iterative method as in Sect.\ \ref{sec:stationary-iterations}. For the chronological history of the development of the CG method, see an excellent review by Golub and O’Leary \BookCite{80}.}{CG 方法 \BookCite{95} 由 Hestenes (1906--1991) 和 Stiefel (1909--1978) 于 1952 年提出, 是最著名的 Krylov 子空间方法. 在精确算术下, 与其他 Krylov 子空间方法一样, CG 方法在有限次迭代内给出精确解. 这意味着 CG 方法具有双重特点: 既是 Section\ \ref{sec:direct-methods} 中的直接方法, 也是 Section\ \ref{sec:stationary-iterations} 中的迭代方法. 关于 CG 方法发展的编年史, 见 Golub 和 O’Leary 的精彩综述 \BookCite{80}.}
% block: cg-lanczos-residual-matrix page:86 type:ordinary-paragraph
\BT{In what follows, the CG method is derived from the Lanczos process (Algorithm\ \ref{alg:lanczos}). From \eqref{eq:lanczos-matrix-relation}, the Lanczos process in matrix form is given by}{下面从 Lanczos 过程 (Algorithm\ \ref{alg:lanczos}) 推导 CG 方法. 由 \eqref{eq:lanczos-matrix-relation}, Lanczos 过程的矩阵形式为}
\begin{equation}\label{eq:cg-lanczos-residual-matrix}AR_n=R_{n+1}T_{n+1,n},\end{equation}
\BT{where}{其中}
\[\begin{gathered}
R_n:=[\bm r_0,\bm r_1,\ldots,\bm r_{n-1}],\\
T_{n+1,n}=\begin{bmatrix}t_{0,0}&t_{0,1}&&\\-\alpha_0^{-1}&t_{1,1}&\ddots&\\&-\alpha_1^{-1}&\ddots&t_{n-2,n-1}\\&&\ddots&t_{n-1,n-1}\\&&&-\alpha_{n-1}^{-1}\end{bmatrix},\\
t_{k-1,k}:=\frac{(\bm r_{k-1},A\bm r_k)}{(\bm r_{k-1},\bm r_{k-1})},\quad(k=1,2,\ldots,n-1),\\
t_{k,k}:=\frac{(\bm r_k,A\bm r_k)}{(\bm r_k,\bm r_k)},\quad(k=0,1,\ldots,n-1).
\end{gathered}\]
\BT{Here, the scaling parameters ($\alpha_k$'s) have not been determined yet. Let $\bm x_k$ be the $k$th approximate solution of $A\bm x=\bm b$ such that $\bm x_k$'s satisfy the following equation:}{这里尺度参数 ($\alpha_k$) 尚未确定. 设 $\bm x_k$ 为 $A\bm x=\bm b$ 的第 $k$ 个近似解, 且这些 $\bm x_k$ 满足下列等式:}
\[R_n=[\bm b-A\bm x_0,\bm b-A\bm x_1,\ldots,\bm b-A\bm x_{n-1}]=\bm b\bm1_n^{\top}-AX_n,\]
\BT{where $\bm1_n:=[1,1,\ldots,1]^{\top}\in\mathbb R^n$ and $X_n:=[\bm x_0,\bm x_1,\ldots,\bm x_{n-1}]$. Notice that $R_n$ is the matrix given in \eqref{eq:cg-lanczos-residual-matrix}. Then, it follows that}{其中 $\bm1_n:=[1,1,\ldots,1]^{\top}\in\mathbb R^n$, $X_n:=[\bm x_0,\bm x_1,\ldots,\bm x_{n-1}]$. 注意, $R_n$ 就是 \eqref{eq:cg-lanczos-residual-matrix} 中给出的矩阵. 于是有}
\begin{equation}\label{eq:cg-residual-solution-matrix}
\begin{aligned}
AR_n&=R_{n+1}T_{n+1,n}=(\bm b\bm1_{n+1}^{\top}-AX_{n+1})T_{n+1,n}\\
\Leftrightarrow R_n&=A^{-1}(\bm b\bm1_{n+1}^{\top}-AX_{n+1})T_{n+1,n}=\bm x\bm1_{n+1}^{\top}T_{n+1,n}-X_{n+1}T_{n+1,n}.
\end{aligned}\end{equation}
% block: cg-scaling-and-ldu pages:86,87 type:ordinary-paragraph
\BT{Now we use the condition $\bm1_{n+1}^{\top}T_{n+1,n}=[0,0,\ldots,0]$ to determine scaling parameters $\alpha_k$, i.e., $\alpha_0^{-1}=t_{0,0}$ and $\alpha_{k-1}^{-1}=t_{k-1,k-1}+t_{k-2,k-1}$ for $k=2,3,\ldots,n$. Then, the condition and \eqref{eq:cg-residual-solution-matrix} yield}{现在用条件 $\bm1_{n+1}^{\top}T_{n+1,n}=[0,0,\ldots,0]$ 确定尺度参数 $\alpha_k$, 即 $\alpha_0^{-1}=t_{0,0}$, 且对 $k=2,3,\ldots,n$, 有 $\alpha_{k-1}^{-1}=t_{k-1,k-1}+t_{k-2,k-1}$. 于是, 这一条件和 \eqref{eq:cg-residual-solution-matrix} 给出}
\begin{equation}\label{eq:cg-residual-negative-product}R_n=-X_{n+1}T_{n+1,n},\end{equation}
\BT{where the matrix $T_{n+1,n}$ in \eqref{eq:cg-lanczos-residual-matrix} is rewritten as}{其中 \eqref{eq:cg-lanczos-residual-matrix} 中的矩阵 $T_{n+1,n}$ 改写为}
\[T_{n+1,n}=\begin{bmatrix}
\alpha_0^{-1}&t_{0,1}&&\\-\alpha_0^{-1}&\alpha_1^{-1}-t_{0,1}&\ddots&\\
&-\alpha_1^{-1}&\ddots&t_{n-2,n-1}\\&&\ddots&\alpha_{n-1}^{-1}-t_{n-2,n-1}\\&&&-\alpha_{n-1}^{-1}
\end{bmatrix}\]
\BT{and the matrix $T_{n+1,n}$ has the following LDU decomposition:}{且矩阵 $T_{n+1,n}$ 具有下列 LDU 分解:}
\begin{equation}\label{eq:cg-lanczos-ldu}
T_{n+1,n}=\underbrace{\begin{bmatrix}1&&\\-1&1&\\\ddots&\ddots&\\&-1&1\end{bmatrix}}_L
\underbrace{\begin{bmatrix}\alpha_0^{-1}&&&\\&\alpha_1^{-1}&&\\&&\ddots&\\&&&\alpha_{n-1}^{-1}\end{bmatrix}}_{D^{-1}}
\underbrace{\begin{bmatrix}1&\alpha_0t_{0,1}&&\\&1&\ddots&\\&&\ddots&\alpha_{n-1}t_{n-1,n}\\&&&1\end{bmatrix}}_U,
\end{equation}
\BT{where $L$ is the $n+1$-by-$n$ lower bidiagonal matrix, $D^{-1}$ is the $n$-by-$n$ diagonal matrix, and $U$ is the $n$-by-$n$ upper bidiagonal matrix. Thus, we obtain}{其中 $L$ 为 $n+1$ 行 $n$ 列的下双对角矩阵, $D^{-1}$ 为 $n$ 行 $n$ 列的对角矩阵, $U$ 为 $n$ 行 $n$ 列的上双对角矩阵. 因此得到}
\begin{equation}\label{eq:cg-residual-ldu}R_n=-X_{n+1}LD^{-1}U.\end{equation}
% block: cg-three-recurrences page:87 type:ordinary-paragraph
\BT{We now define $P_n(=[\bm p_0,\bm p_1,\ldots,\bm p_{n-1}]):=R_nU^{-1}$. Then, \eqref{eq:cg-residual-ldu} can be rewritten as}{现在定义 $P_n(=[\bm p_0,\bm p_1,\ldots,\bm p_{n-1}]):=R_nU^{-1}$. 则 \eqref{eq:cg-residual-ldu} 可改写为}
\begin{equation}\label{eq:cg-direction-matrix}P_nD=X_{n+1}(-L),\end{equation}
\BT{which is equivalent to}{它等价于}
\begin{equation}\label{eq:cg-solution-recurrence}\bm x_{k+1}=\bm x_k+\alpha_k\bm p_k,\quad k=0,1,\ldots,n-1.\end{equation}
\BT{The relation between $\bm r_k=\bm b-A\bm x_k$ and \eqref{eq:cg-solution-recurrence} yields}{由关系 $\bm r_k=\bm b-A\bm x_k$ 和 \eqref{eq:cg-solution-recurrence} 得}
\begin{equation}\label{eq:cg-residual-recurrence}\bm r_{k+1}=\bm r_k-\alpha_kA\bm p_k,\quad k=0,1,\ldots,n-1.\end{equation}
\BT{The definition $P_n=R_nU^{-1}$ leads to $P_nU=R_n$, which is equivalent to}{定义 $P_n=R_nU^{-1}$ 给出 $P_nU=R_n$, 它等价于}
\begin{equation}\label{eq:cg-initial-direction}\bm p_0=\bm r_0,\end{equation}
\begin{equation}\label{eq:cg-direction-recurrence}\bm p_k=\bm r_k+\beta_{k-1}\bm p_{k-1},\quad k=1,2,\ldots,n-1,\end{equation}
\BT{where $\beta_{k-1}:=-\alpha_{k-1}t_{k-1,k}$. As described above, computing $\alpha_k$ requires $A\bm r_k$ because $\alpha_k=(t_{k,k}+t_{k-1,k})^{-1}$, which means two matrix--vector multiplications are required per iteration step, i.e., $A\bm r_k$ and $A\bm p_k$ if we use the iterates \eqref{eq:cg-solution-recurrence}--\eqref{eq:cg-direction-recurrence}. In what follows, we will see that $A\bm r_k$ is not required. To see this, first we derive a computational formula for $\alpha_k$.}{其中 $\beta_{k-1}:=-\alpha_{k-1}t_{k-1,k}$. 如上所述, 计算 $\alpha_k$ 需要 $A\bm r_k$, 因为 $\alpha_k=(t_{k,k}+t_{k-1,k})^{-1}$; 这意味着若使用迭代式 \eqref{eq:cg-solution-recurrence}--\eqref{eq:cg-direction-recurrence}, 每次迭代需要两次矩阵向量乘法, 即 $A\bm r_k$ 和 $A\bm p_k$. 下面将说明不需要 $A\bm r_k$. 为此先推导 $\alpha_k$ 的计算公式.}
% block: cg-step-length page:88 type:ordinary-paragraph
\BT{From \eqref{eq:cg-lanczos-residual-matrix}, \eqref{eq:cg-lanczos-ldu}, and the definition of $P_n$, it follows that}{由 \eqref{eq:cg-lanczos-residual-matrix}, \eqref{eq:cg-lanczos-ldu} 及 $P_n$ 的定义, 有}
\[AP_n=R_{n+1}LD^{-1}.\]
\BT{This leads to}{由此得到}
\begin{equation}\label{eq:cg-projected-diagonal-product}P_n^{\mathrm H}AP_nD=U^{-\mathrm H}R_n^{\mathrm H}R_{n+1}L.\end{equation}
\BT{Since $R_n=[\bm r_0,\bm r_1,\ldots,\bm r_{n-1}]$ in \eqref{eq:cg-lanczos-residual-matrix} is generated by the Lanczos process, we have $\bm r_i^{\mathrm H}\bm r_j=0$ for $i\ne j$. In matrix form, this means $(i,j)$ elements of $R_n^{\mathrm H}R_{n+1}$ is zero for $i\ne j$. This fact yields that the $(k+1,k+1)$ element of $U^{-\mathrm H}R_n^{\mathrm H}R_{n+1}L$ in \eqref{eq:cg-projected-diagonal-product} is $\bm r_k^{\mathrm H}\bm r_k$ because $L$ and $U^{-\mathrm H}$ are the lower triangular matrices with all the diagonals being one. Thus from the $(k+1,k+1)$ element in \eqref{eq:cg-projected-diagonal-product}, we have $\bm p_k^{\mathrm H}A\bm p_k\alpha_k=\bm r_k^{\mathrm H}\bm r_k$, or}{由于 \eqref{eq:cg-lanczos-residual-matrix} 中的 $R_n=[\bm r_0,\bm r_1,\ldots,\bm r_{n-1}]$ 由 Lanczos 过程生成, 当 $i\ne j$ 时有 $\bm r_i^{\mathrm H}\bm r_j=0$. 用矩阵形式表示, 即当 $i\ne j$ 时 $R_n^{\mathrm H}R_{n+1}$ 的 $(i,j)$ 元素为零. 由于 $L$ 和 $U^{-\mathrm H}$ 都是对角元素全为一的下三角矩阵, 这一事实给出 \eqref{eq:cg-projected-diagonal-product} 中 $U^{-\mathrm H}R_n^{\mathrm H}R_{n+1}L$ 的 $(k+1,k+1)$ 元素为 $\bm r_k^{\mathrm H}\bm r_k$. 因此由 \eqref{eq:cg-projected-diagonal-product} 的 $(k+1,k+1)$ 元素有 $\bm p_k^{\mathrm H}A\bm p_k\alpha_k=\bm r_k^{\mathrm H}\bm r_k$, 即}
\begin{equation}\label{eq:cg-alpha}\alpha_k=\frac{\bm r_k^{\mathrm H}\bm r_k}{\bm p_k^{\mathrm H}A\bm p_k}.\end{equation}
\BT{Since $A$ is Hermitian, $\alpha_k$ is a real number for all $k$, which will be later used for deriving the computational formula for $\beta_k$.}{由于 $A$ 为 Hermitian 矩阵, 对所有 $k$, $\alpha_k$ 都是实数, 稍后推导 $\beta_k$ 的计算公式时将使用这一点.}
% block: cg-direction-coefficient page:88 type:ordinary-paragraph
\BT{Next, we give a computational formula for $\beta_k$. From \eqref{eq:cg-lanczos-residual-matrix}, it follows that}{接着给出 $\beta_k$ 的计算公式. 由 \eqref{eq:cg-lanczos-residual-matrix} 得}
\begin{equation}\label{eq:cg-residual-projected-matrix}R_n^{\mathrm H}AR_n=R_n^{\mathrm H}R_{n+1}T_{n+1,n},\end{equation}
\BT{The $(k+1,k)$ element of $R_n^{\mathrm H}AR_n$ is $\bm r_k^{\mathrm H}A\bm r_{k-1}$. As seen in \eqref{eq:cg-projected-diagonal-product}, $R_n^{\mathrm H}R_{n+1}$ is a matrix whose $(k,k)$ elements are $\bm r_{k-1}^{\mathrm H}\bm r_{k-1}$ for $k=1,2,\ldots,n$ and the other elements are zero. From this, it is seen that the $(k+1,k)$ element of $R_n^{\mathrm H}R_{n+1}T_{n+1,n}$ is $\bm r_k^{\mathrm H}\bm r_k\times(-\alpha_{k-1})^{-1}$. Thus we obtain $\bm r_k^{\mathrm H}A\bm r_{k-1}=\bm r_k^{\mathrm H}\bm r_k\times(-\alpha_{k-1})^{-1}$. Because $\alpha_{k-1}$ is a real number as mentioned before, $\bm r_k^{\mathrm H}A\bm r_{k-1}$ is also a real number. Thus $\bm r_k^{\mathrm H}A\bm r_{k-1}=(\bm r_k^{\mathrm H}A\bm r_{k-1})^{\mathrm H}=\bm r_{k-1}^{\mathrm H}A^{\mathrm H}\bm r_k=\bm r_{k-1}^{\mathrm H}A\bm r_k$. From this fact it follows that $\bm r_{k-1}^{\mathrm H}A\bm r_k=\bm r_k^{\mathrm H}\bm r_k\times(-\alpha_{k-1})^{-1}$. Recalling $\beta_{k-1}=-\alpha_{k-1}t_{k-1,k}$ yields}{$R_n^{\mathrm H}AR_n$ 的 $(k+1,k)$ 元素为 $\bm r_k^{\mathrm H}A\bm r_{k-1}$. 如 \eqref{eq:cg-projected-diagonal-product} 所示, $R_n^{\mathrm H}R_{n+1}$ 在 $k=1,2,\ldots,n$ 时的 $(k,k)$ 元素为 $\bm r_{k-1}^{\mathrm H}\bm r_{k-1}$, 其余元素为零. 由此可见, $R_n^{\mathrm H}R_{n+1}T_{n+1,n}$ 的 $(k+1,k)$ 元素为 $\bm r_k^{\mathrm H}\bm r_k\times(-\alpha_{k-1})^{-1}$. 因此得到 $\bm r_k^{\mathrm H}A\bm r_{k-1}=\bm r_k^{\mathrm H}\bm r_k\times(-\alpha_{k-1})^{-1}$. 由于前述 $\alpha_{k-1}$ 为实数, $\bm r_k^{\mathrm H}A\bm r_{k-1}$ 也是实数. 因而 $\bm r_k^{\mathrm H}A\bm r_{k-1}=(\bm r_k^{\mathrm H}A\bm r_{k-1})^{\mathrm H}=\bm r_{k-1}^{\mathrm H}A^{\mathrm H}\bm r_k=\bm r_{k-1}^{\mathrm H}A\bm r_k$. 由此得 $\bm r_{k-1}^{\mathrm H}A\bm r_k=\bm r_k^{\mathrm H}\bm r_k\times(-\alpha_{k-1})^{-1}$. 回顾 $\beta_{k-1}=-\alpha_{k-1}t_{k-1,k}$, 得}
\begin{equation}\label{eq:cg-beta}
\beta_{k-1}=-\alpha_{k-1}t_{k-1,k}=\frac{\bm r_k^{\mathrm H}\bm r_k}{\bm r_{k-1}^{\mathrm H}A\bm r_k}\frac{\bm r_{k-1}^{\mathrm H}A\bm r_k}{\bm r_{k-1}^{\mathrm H}\bm r_{k-1}}=\frac{\bm r_k^{\mathrm H}\bm r_k}{\bm r_{k-1}^{\mathrm H}\bm r_{k-1}}.
\end{equation}
\BT{From \eqref{eq:cg-solution-recurrence}--\eqref{eq:cg-direction-recurrence}, \eqref{eq:cg-alpha}, and \eqref{eq:cg-beta}, the algorithm of the CG method is obtained, which is listed in Algorithm\ \ref{alg:cg}. Properties of the CG method are described in Proposition\ \ref{prop:cg-orthogonality}.}{由 \eqref{eq:cg-solution-recurrence}--\eqref{eq:cg-direction-recurrence}, \eqref{eq:cg-alpha} 和 \eqref{eq:cg-beta} 得到 CG 方法的算法, 见 Algorithm\ \ref{alg:cg}. CG 方法的性质见 Proposition\ \ref{prop:cg-orthogonality}.}
% block: cg-algorithm page:89 type:ordinary-paragraph
\begin{algorithm}[H]\caption{\BT{The CG method}{CG 方法}}\label{alg:cg}
\begin{algorithmic}[1]
\Statex\textbf{Input:} $\bm x_0\in\mathbb C^N,\ \beta_{-1}=0,\ \bm p_{-1}=\bm0,\ \bm r_0=\bm b-A\bm x_0$
\Statex\textbf{Output:} $\bm x_n$
\For{$n=0,1,2,\ldots$, \BT{until convergence}{直至收敛}}
\State $\bm p_n=\bm r_n+\beta_{n-1}\bm p_{n-1}$
\State $\alpha_n=\frac{(\bm r_n,\bm r_n)}{(\bm p_n,A\bm p_n)}$
\State $\bm x_{n+1}=\bm x_n+\alpha_n\bm p_n$
\State $\bm r_{n+1}=\bm r_n-\alpha_nA\bm p_n$
\State $\beta_n=\frac{(\bm r_{n+1},\bm r_{n+1})}{(\bm r_n,\bm r_n)}$
\EndFor\end{algorithmic}\end{algorithm}
% block: cg-orthogonality-proposition page:88 type:proposition
\begin{Proposition}\label{prop:cg-orthogonality}
\BT{Let $\bm p_n,\bm r_n$ be the vectors in the CG method, then:}{设 $\bm p_n,\bm r_n$ 为 CG 方法中的向量, 则:}
\begin{enumerate}\item $(\bm r_i,\bm r_j)=0\quad\text{for }i\ne j$,\item $(\bm p_i,A\bm p_j)=0\quad\text{for }i\ne j$.\end{enumerate}
\end{Proposition}
% block: cg-orthogonality-proof pages:88,89 type:proof
\begin{Proof}
\BT{The first property readily follows from the fact that \eqref{eq:cg-lanczos-residual-matrix} is obtained from the Lanczos process. For the second property, it follows from \eqref{eq:cg-projected-diagonal-product} that we have}{第一个性质由 \eqref{eq:cg-lanczos-residual-matrix} 来自 Lanczos 过程这一事实立即得到. 对第二个性质, 由 \eqref{eq:cg-projected-diagonal-product} 有}
\[P_n^{\mathrm H}AP_n=U^{-\mathrm H}R_n^{\mathrm H}R_{n+1}LD^{-1}.\]
\BT{$P_n^{\mathrm H}AP_n$ is Hermitian because $A$ is Hermitian. Thus, $U^{-\mathrm H}R_n^{\mathrm H}R_{n+1}LD^{-1}$ is Hermitian. Since the $(i,j)$ element $R_n^{\mathrm H}R_{n+1}$ is zero for $i\ne j$, matrix $U^{-\mathrm H}R_n^{\mathrm H}R_{n+1}LD^{-1}$ is a lower triangular matrix. This means that $P_n^{\mathrm H}AP_n$ is a Hermitian and lower triangular matrix. Thus $P_n^{\mathrm H}AP_n$ is a diagonal matrix, which is equivalent to the second property.}{由于 $A$ 为 Hermitian 矩阵, $P_n^{\mathrm H}AP_n$ 为 Hermitian 矩阵. 因此 $U^{-\mathrm H}R_n^{\mathrm H}R_{n+1}LD^{-1}$ 为 Hermitian 矩阵. 由于 $i\ne j$ 时 $R_n^{\mathrm H}R_{n+1}$ 的 $(i,j)$ 元素为零, 矩阵 $U^{-\mathrm H}R_n^{\mathrm H}R_{n+1}LD^{-1}$ 是下三角矩阵. 这表明 $P_n^{\mathrm H}AP_n$ 既是 Hermitian 矩阵, 又是下三角矩阵. 因而 $P_n^{\mathrm H}AP_n$ 为对角矩阵, 这等价于第二个性质.}
\end{Proof}
% block: cg-krylov-corollary page:89 type:corollary
\BT{From Proposition\ \ref{prop:cg-orthogonality}, we have the following properties of the CG method:}{由 Proposition\ \ref{prop:cg-orthogonality}, 有 CG 方法的下列性质:}
\begin{Corollary}\label{cor:cg-krylov-orthogonality}
\BT{Let $\bm p_n$ and $\bm r_n$ be the vectors in the CG method, then:}{设 $\bm p_n$ 和 $\bm r_n$ 为 CG 方法中的向量, 则:}
\begin{enumerate}\item $\bm r_n\perp\mathcal K_n(A,\bm r_0)$,\item $A\bm p_n\perp\mathcal K_n(A,\bm r_0)$.\end{enumerate}
\end{Corollary}
% block: cg-krylov-corollary-proof page:89 type:proof
\begin{Proof}
\BT{From \eqref{eq:cg-residual-recurrence} and \eqref{eq:cg-initial-direction}, we have $\bm r_1=\bm r_0-\alpha_0A\bm r_0$. This means that $\operatorname{span}\{\bm r_0,\bm r_1\}=\mathcal K_2(A,\bm r_0)$. Similarly, $\operatorname{span}\{\bm r_0,\bm r_1,\ldots,\bm r_{n-1}\}=\mathcal K_n(A,\bm r_0)$. From the first property of Proposition\ \ref{prop:cg-orthogonality}, it follows that $\bm r_n\perp\bm r_0,\bm r_1,\ldots,\bm r_{n-1}$. Thus $\bm r_n\perp\mathcal K_n(A,\bm r_0)$. Next, from \eqref{eq:cg-residual-recurrence} and the first property of Proposition\ \ref{prop:cg-orthogonality}, it follows that $A\bm p_n=(\bm r_n-\bm r_{n+1})/\alpha_n\perp\bm r_0,\bm r_1,\ldots,\bm r_{n-1}$, and thus $A\bm p_n\perp\mathcal K_n(A,\bm r_0)$.}{由 \eqref{eq:cg-residual-recurrence} 和 \eqref{eq:cg-initial-direction} 有 $\bm r_1=\bm r_0-\alpha_0A\bm r_0$. 这意味着 $\operatorname{span}\{\bm r_0,\bm r_1\}=\mathcal K_2(A,\bm r_0)$. 类似地, $\operatorname{span}\{\bm r_0,\bm r_1,\ldots,\bm r_{n-1}\}=\mathcal K_n(A,\bm r_0)$. 由 Proposition\ \ref{prop:cg-orthogonality} 的第一个性质, 得 $\bm r_n\perp\bm r_0,\bm r_1,\ldots,\bm r_{n-1}$. 因此 $\bm r_n\perp\mathcal K_n(A,\bm r_0)$. 接着, 由 \eqref{eq:cg-residual-recurrence} 和 Proposition\ \ref{prop:cg-orthogonality} 的第一个性质, 得 $A\bm p_n=(\bm r_n-\bm r_{n+1})/\alpha_n\perp\bm r_0,\bm r_1,\ldots,\bm r_{n-1}$, 因而 $A\bm p_n\perp\mathcal K_n(A,\bm r_0)$.}
\end{Proof}
% block: cg-search-space-and-optimality page:89 type:ordinary-paragraph
\BT{From the proof of Corollary\ \ref{cor:cg-krylov-orthogonality}, it is easy to see that}{由 Corollary\ \ref{cor:cg-krylov-orthogonality} 的证明容易看出}
\begin{equation}\label{eq:cg-krylov-membership}
\bm r_n\in\mathcal K_{n+1}(A,\bm r_0),\quad\bm p_n\in\mathcal K_{n+1}(A,\bm r_0),\quad\bm x_n-\bm x_0\in\mathcal K_n(A,\bm r_0).
\end{equation}
\BT{If the coefficient matrix $A$ is a Hermitian (or real symmetric) positive definite matrix, then the CG method produces the optimal approximate solution in terms of $A$-norm of the error as described below.}{若系数矩阵 $A$ 为 Hermitian (或实对称) 正定矩阵, 则 CG 方法给出关于误差 $A$-范数的最优近似解, 如下所述.}
% block: cg-optimality-theorem page:89 type:theorem
\begin{Theorem}\label{thm:cg-energy-optimality}
\BT{Let $A\bm x=\bm b$ be Hermitian positive definite linear systems. Then the CG method produces approximate solutions such that $A$-norm of the error is minimized:}{设 $A\bm x=\bm b$ 为 Hermitian 正定线性方程组. 则 CG 方法产生使误差的 $A$-范数达到极小值的近似解:}
\[\min_{\bm x_n\in\bm x_0+\mathcal K_n(A,\bm r_0)}\|\bm e_n\|_A,\]
\BT{where $\|\bm e_n\|_A=(\bm e_n^{\mathrm H}A\bm e_n)^{1/2}$ and $\bm e_n=\bm x-\bm x_n$ is the error of the $n$th approximate solution of the CG method.}{其中 $\|\bm e_n\|_A=(\bm e_n^{\mathrm H}A\bm e_n)^{1/2}$, $\bm e_n=\bm x-\bm x_n$ 为 CG 方法第 $n$ 个近似解的误差.}
\end{Theorem}
% block: cg-optimality-proof page:90 type:proof
\begin{Proof}
\BT{Since $A$ is Hermitian positive definite, the Cholesky decomposition exists, i.e., $A=LL^{\mathrm H}$. (The Cholesky decomposition for a real symmetric positive definite matrix is described in Sect.\ \ref{subsec:cholesky-decomposition}.) Recall that the approximate solutions of the CG method satisfy}{由于 $A$ 为 Hermitian 正定矩阵, Cholesky 分解存在, 即 $A=LL^{\mathrm H}$. (实对称正定矩阵的 Cholesky 分解见 Section\ \ref{subsec:cholesky-decomposition}.) 回顾 CG 方法的近似解满足}
\[\bm x_n=\bm x_0+\bm z_n,\quad\bm z_n\in\mathcal K_n(A,\bm r_0),\]
\BT{where $\bm z_n$ is determined by $\bm r_n\perp\mathcal K_n(A,\bm r_0)$, see Corollary\ \ref{cor:cg-krylov-orthogonality}. Let $V_n$ be a matrix whose column vectors correspond to orthonormalized basis vectors of $\mathcal K_n(A,\bm r_0)$. Then,}{其中 $\bm z_n$ 由 $\bm r_n\perp\mathcal K_n(A,\bm r_0)$ 确定, 见 Corollary\ \ref{cor:cg-krylov-orthogonality}. 设 $V_n$ 为列向量对应于 $\mathcal K_n(A,\bm r_0)$ 的标准正交基向量的矩阵. 则}
\begin{equation}\label{eq:cg-orthonormal-coordinate-solution}\bm x_n=\bm x_0+V_n\bm y_n,\quad\bm y_n\in\mathbb C^n,\end{equation}
\BT{where $\bm y_n$ is determined by $V_n^{\mathrm H}\bm r_n=\bm0$. From \eqref{eq:cg-orthonormal-coordinate-solution} and $\bm e_n=\bm x-\bm x_n$, it follows that}{其中 $\bm y_n$ 由 $V_n^{\mathrm H}\bm r_n=\bm0$ 确定. 由 \eqref{eq:cg-orthonormal-coordinate-solution} 和 $\bm e_n=\bm x-\bm x_n$, 得}
\[\bm e_n=\bm e_0-V_n\bm y_n,\]
\BT{where $\bm y_n$ is determined by $V_n^{\mathrm H}LL^{\mathrm H}\bm e_n=\bm0$ since $V_n^{\mathrm H}\bm r_n=\bm0$ and $\bm r_n=\bm b-A\bm x_n=A\bm x-A\bm x_n=A\bm e_n=LL^{\mathrm H}\bm e_n$. Here we used the Cholesky decomposition $A=LL^{\mathrm H}$. Equivalently, we have}{其中 $\bm y_n$ 由 $V_n^{\mathrm H}LL^{\mathrm H}\bm e_n=\bm0$ 确定, 因为 $V_n^{\mathrm H}\bm r_n=\bm0$, 且 $\bm r_n=\bm b-A\bm x_n=A\bm x-A\bm x_n=A\bm e_n=LL^{\mathrm H}\bm e_n$. 这里使用了 Cholesky 分解 $A=LL^{\mathrm H}$. 等价地, 有}
\[L^{\mathrm H}\bm e_n=L^{\mathrm H}\bm e_0-L^{\mathrm H}V_n\bm y_n,\]
\BT{where $\bm y_n$ is determined by $(L^{\mathrm H}V_n)^{\mathrm H}L^{\mathrm H}\bm e_n=\bm0$. This means that the CG method finds approximate solutions such that}{其中 $\bm y_n$ 由 $(L^{\mathrm H}V_n)^{\mathrm H}L^{\mathrm H}\bm e_n=\bm0$ 确定. 这意味着 CG 方法寻找满足下列条件的近似解}
\begin{equation}\label{eq:cg-transformed-least-squares}\min_{\bm y_n\in\mathbb C^n}\|L^{\mathrm H}\bm e_n\|,\end{equation}
\BT{because $(L^{\mathrm H}V_n)^{\mathrm H}L^{\mathrm H}\bm e_n=\bm0$ is equivalent to the normal equation $(L^{\mathrm H}V_n)^{\mathrm H}(L^{\mathrm H}V_n)\bm y_n=(L^{\mathrm H}V_n)^{\mathrm H}L^{\mathrm H}\bm e_0$. The minimization \eqref{eq:cg-transformed-least-squares} is equivalent to}{因为 $(L^{\mathrm H}V_n)^{\mathrm H}L^{\mathrm H}\bm e_n=\bm0$ 等价于正规方程 $(L^{\mathrm H}V_n)^{\mathrm H}(L^{\mathrm H}V_n)\bm y_n=(L^{\mathrm H}V_n)^{\mathrm H}L^{\mathrm H}\bm e_0$. 极小化问题 \eqref{eq:cg-transformed-least-squares} 等价于}
\[\min_{\bm x_n\in\bm x_0+\mathcal K_n(A,\bm r_0)}\|L^{\mathrm H}\bm e_n\|.\]
\BT{Since $\|L^{\mathrm H}\bm e_n\|^2=(L^{\mathrm H}\bm e_n)^{\mathrm H}(L^{\mathrm H}\bm e_n)=\bm e_n^{\mathrm H}LL^{\mathrm H}\bm e_n=\bm e_n^{\mathrm H}A\bm e_n$, we finally have}{由于 $\|L^{\mathrm H}\bm e_n\|^2=(L^{\mathrm H}\bm e_n)^{\mathrm H}(L^{\mathrm H}\bm e_n)=\bm e_n^{\mathrm H}LL^{\mathrm H}\bm e_n=\bm e_n^{\mathrm H}A\bm e_n$, 最终有}
\[\min_{\bm x_n\in\bm x_0+\mathcal K_n(A,\bm r_0)}\|\bm e_n\|_A,\]
\BT{which concludes the proof.}{证明完毕.}\end{Proof}
% block: cg-polynomial-representation pages:90,91 type:ordinary-paragraph
\BT{Theorem\ \ref{thm:cg-energy-optimality} shows an optimality of the CG method. On the other hand, we cannot see how fast the CG method converges. To see the convergence rate later, the matrix polynomial representation of the CG method is described. The residual vector $\bm r_n$ and the auxiliary vector $\bm p_n$ of the CG method can be expressed by using two polynomials $R_n$ and $P_n$,}{Theorem\ \ref{thm:cg-energy-optimality} 给出 CG 方法的一种最优性. 另一方面, 尚无法由此看出 CG 方法收敛得多快. 为稍后考察收敛速度, 先介绍 CG 方法的矩阵多项式表示. CG 方法的残差向量 $\bm r_n$ 和辅助向量 $\bm p_n$ 可用两个多项式 $R_n$ 和 $P_n$ 表示,}
\begin{equation}\label{eq:cg-residual-direction-polynomials}\bm r_n=R_n(A)\bm r_0,\quad\bm p_n=P_n(A)\bm r_0,\end{equation}
\BT{where $R_n(\lambda)$ denotes Lanczos polynomials, which satisfy the following three-term recurrence relations:}{其中 $R_n(\lambda)$ 表示 Lanczos 多项式, 它满足下列三项递推关系:}
\begin{equation}\label{eq:cg-polynomial-initial-zero}R_0(\lambda)=1,\end{equation}
\begin{equation}\label{eq:cg-polynomial-initial-one}R_1(\lambda)=(1-\alpha_0\lambda)R_0(\lambda),\end{equation}
\begin{equation}\label{eq:cg-polynomial-three-term}
\begin{aligned}R_n(\lambda)&=\left(1+\frac{\beta_{n-2}}{\alpha_{n-2}}\alpha_{n-1}-\alpha_{n-1}\lambda\right)R_{n-1}(\lambda)\\
&\quad-\frac{\beta_{n-2}}{\alpha_{n-2}}\alpha_{n-1}R_{n-2}(\lambda),\quad n=2,3,\ldots,
\end{aligned}\end{equation}
\BT{where $R_n(0)=1$. $R_n$ and $P_n$ satisfy the following two-term recurrences:}{其中 $R_n(0)=1$. $R_n$ 和 $P_n$ 满足下列两项递推关系:}
\begin{equation}\label{eq:cg-two-term-polynomial-initial}R_0(\lambda)=1,\quad P_0(\lambda)=1,\end{equation}
\begin{equation}\label{eq:cg-residual-polynomial-recurrence}R_n(\lambda)=R_{n-1}(\lambda)-\alpha_{n-1}\lambda P_{n-1}(\lambda),\end{equation}
\begin{equation}\label{eq:cg-direction-polynomial-recurrence}P_n(\lambda)=R_n(\lambda)+\beta_{n-1}P_{n-1}(\lambda)\quad n=1,2,\ldots.\end{equation}
% block: cg-convergence-theorem page:91 type:theorem
\BT{The rate of the convergence of the CG method is described in Theorem\ \ref{thm:cg-condition-number-bound}.}{CG 方法的收敛速度见 Theorem\ \ref{thm:cg-condition-number-bound}.}
\Needspace{7\baselineskip}
\begin{Theorem}\label{thm:cg-condition-number-bound}
\BT{The $A$-norm of the error of the CG method satisfies}{CG 方法误差的 $A$-范数满足}
\begin{equation}\label{eq:cg-condition-number-bound}
\|\bm e_n\|_A\le2\left(\frac{\sqrt\kappa-1}{\sqrt\kappa+1}\right)^n\|\bm e_0\|_A,
\end{equation}
\BT{where $\kappa$ is the condition number of $A$.}{其中 $\kappa$ 为 $A$ 的条件数.}
\end{Theorem}
% block: cg-convergence-proof-diagonalization page:91 type:proof
\begin{Proof}
\BT{Since the residual vector of the CG method can be described as $\bm r_n=R(A)\bm r_0$ and $\bm r_n=A\bm e_n$, we obtain $\bm e_n=R_n(A)\bm e_0$. From \eqref{eq:cg-condition-number-bound} it follows that}{由于 CG 方法的残差向量可表示为 $\bm r_n=R(A)\bm r_0$, 且 $\bm r_n=A\bm e_n$, 得到 $\bm e_n=R_n(A)\bm e_0$. 由 \eqref{eq:cg-condition-number-bound} 有}
\[\|\bm e_n\|_A=\|R_n(A)\bm e_0\|_A=\min_{p_n\in\mathcal P_n}\|p_n(A)\bm e_0\|_A,\]
\BT{where $\mathcal P_n$ is the set of polynomials of degree $n$ with $\mathcal P_0=1$. Using the diagonalization of Hermitian positive definite matrix $A$, we have $A=PDP^{\mathrm H}$ with unitary matrix $P$ and diagonal matrix $D$ whose diagonal elements are eigenvalues of $A$. Let $A^{1/2}=PD^{1/2}P^{\mathrm H}$, where $D^{1/2}$ is a diagonal matrix whose diagonal elements are square roots of eigenvalues of $A$. Note that all the eigenvalues of $A$ are positive and $A^{1/2}$ is Hermitian, i.e., $(A^{1/2})^{\mathrm H}=A^{1/2}$. We see that $A^{1/2}A^{1/2}=PD^{1/2}P^{\mathrm H}PD^{1/2}P^{\mathrm H}=PD^{1/2}D^{1/2}P^{\mathrm H}=PDP^{\mathrm H}=A$. Then we have}{其中 $\mathcal P_n$ 为次数为 $n$ 且满足 $\mathcal P_0=1$ 的多项式的集合. 利用 Hermitian 正定矩阵 $A$ 的对角化, 有 $A=PDP^{\mathrm H}$, 其中 $P$ 为酉矩阵, $D$ 为对角元素是 $A$ 的特征值的对角矩阵. 令 $A^{1/2}=PD^{1/2}P^{\mathrm H}$, 其中 $D^{1/2}$ 的对角元素为 $A$ 的特征值的平方根. 注意, $A$ 的所有特征值均为正, 且 $A^{1/2}$ 为 Hermitian 矩阵, 即 $(A^{1/2})^{\mathrm H}=A^{1/2}$. 可见 $A^{1/2}A^{1/2}=PD^{1/2}P^{\mathrm H}PD^{1/2}P^{\mathrm H}=PD^{1/2}D^{1/2}P^{\mathrm H}=PDP^{\mathrm H}=A$. 于是有}
% block: cg-convergence-proof-norm-chain page:92 type:proof
\[\begin{aligned}
\|\bm e_n\|_A&=\min_{p_n\in\mathcal P_n}\|p_n(A)\bm e_0\|_A,\\
&=\min_{p_n\in\mathcal P_n}\sqrt{(p_n(A)\bm e_0)^{\mathrm H}A(p_n(A)\bm e_0)}\\
&=\min_{p_n\in\mathcal P_n}\sqrt{(p_n(A)\bm e_0)^{\mathrm H}(A^{1/2})^{\mathrm H}A^{1/2}(p_n(A)\bm e_0)}\\
&=\min_{p_n\in\mathcal P_n}\sqrt{(A^{1/2}p_n(A)\bm e_0)^{\mathrm H}(A^{1/2}p_n(A)\bm e_0)}\\
&=\min_{p_n\in\mathcal P_n}\|A^{1/2}p_n(A)\bm e_0\|.
\end{aligned}\]
\BT{Using $A=PDP^{\mathrm H}$ and $A^{1/2}=PD^{1/2}P^{\mathrm H}$ yields}{利用 $A=PDP^{\mathrm H}$ 和 $A^{1/2}=PD^{1/2}P^{\mathrm H}$, 得}
\[\begin{aligned}
\|\bm e_n\|_A&=\min_{p_n\in\mathcal P_n}\|A^{1/2}p_n(A)\bm e_0\|=\min_{p_n\in\mathcal P_n}\|PD^{1/2}P^{\mathrm H}p_n(PDP^{\mathrm H})\bm e_0\|\\
&=\min_{p_n\in\mathcal P_n}\|PD^{1/2}P^{\mathrm H}Pp_n(D)P^{\mathrm H}\bm e_0\|
=\min_{p_n\in\mathcal P_n}\|PD^{1/2}p_n(D)P^{\mathrm H}\bm e_0\|\\
&=\min_{p_n\in\mathcal P_n}\|Pp_n(D)D^{1/2}P^{\mathrm H}\bm e_0\|
=\min_{p_n\in\mathcal P_n}\|Pp_n(D)P^{\mathrm H}PD^{1/2}P^{\mathrm H}\bm e_0\|\\
&=\min_{p_n\in\mathcal P_n}\|Pp_n(D)P^{\mathrm H}A^{1/2}\bm e_0\|\\
&\le\min_{p_n\in\mathcal P_n}\|Pp_n(D)P^{\mathrm H}\|\|A^{1/2}\bm e_0\|\\
&=\min_{p_n\in\mathcal P_n}\|p_n(D)\|\|\bm e_0\|_A\\
&=\min_{p_n\in\mathcal P_n}\max_{i=1,2,\ldots,N}|p_n(\lambda_i)|\times\|\bm e_0\|_A.
\end{aligned}\]
\BT{Thus, we have}{因此有}
\[\frac{\|\bm e_n\|_A}{\|\bm e_0\|_A}\le\min_{p_n\in\mathcal P_n}\max_{i=1,2,\ldots,N}|p_n(\lambda_i)|.\]
\BT{Let $a$ be the smallest eigenvalue and $b$ be the largest eigenvalue of $A$. Then all the eigenvalues $\lambda_i$ belong to the closed interval $[a,b]$. Thus we have the following inequalities:}{设 $a$ 和 $b$ 分别为 $A$ 的最小和最大特征值. 则所有特征值 $\lambda_i$ 都属于闭区间 $[a,b]$. 因此有下列不等式:}
\[\frac{\|\bm e_n\|_A}{\|\bm e_0\|_A}\le\min_{p_n\in\mathcal P_n}\max_{i=1,2,\ldots,N}|p_n(\lambda_i)|\le\min_{p_n\in\mathcal P_n}\max_{\lambda\in[a,b]}|p_n(\lambda)|.\]
\BT{The solution of the min-max problem can be rewritten using Chebyshev polynomials as follows:}{极小极大问题的解可用 Chebyshev 多项式改写如下:}
\[\min_{p_n\in\mathcal P_n}\max_{\lambda\in[a,b]}|p_n(\lambda)|=\max_{\lambda\in[a,b]}\frac{|T_n(\frac{2\lambda-b-a}{b-a})|}{T_n(\frac{-b-a}{b-a})},\]
% block: cg-convergence-proof-chebyshev page:93 type:proof
\BT{where}{其中}
\begin{equation}\label{eq:chebyshev-explicit-form}
T_n(x)=\frac12\left[\left(x+\sqrt{x^2-1}\right)^n+\left(x-\sqrt{x^2-1}\right)^n\right]\quad\text{for }|x|\ge1.
\end{equation}
\BT{From the property of Chebyshev polynomials, $|T_n(\frac{2\lambda-b-a}{b-a})|\le1$, from which we have}{由 Chebyshev 多项式的性质, $|T_n(\frac{2\lambda-b-a}{b-a})|\le1$, 从而有}
\[\min_{p_n\in\mathcal P_n}\max_{\lambda\in[a,b]}|p_n(\lambda)|\le\frac1{|T_n(\frac{-b-a}{b-a})|}=\frac1{|T_n(\frac{-\kappa-1}{\kappa-1})|},\]
\BT{where $\kappa=b/a$ is the condition number of $A$. Using \eqref{eq:chebyshev-explicit-form} yields}{其中 $\kappa=b/a$ 是 $A$ 的条件数. 利用 \eqref{eq:chebyshev-explicit-form}, 得}
\[\begin{aligned}
\min_{p_n\in\mathcal P_n}\max_{\lambda\in[a,b]}|p_n(\lambda)|&\le\frac1{|T_n(\frac{-\kappa-1}{\kappa-1})|}
=\frac2{\left(\frac{\sqrt\kappa-1}{\sqrt\kappa+1}\right)^n+\left(\frac{\sqrt\kappa+1}{\sqrt\kappa-1}\right)^n}
\le\frac2{\left(\frac{\sqrt\kappa+1}{\sqrt\kappa-1}\right)^n}\\
&=2\left(\frac{\sqrt\kappa-1}{\sqrt\kappa+1}\right)^n,
\end{aligned}\]
\BT{which concludes the proof.}{证明完毕.}\end{Proof}
% block: cg-preconditioning-transformation page:93 type:ordinary-paragraph
\BT{Theorem\ \ref{thm:cg-condition-number-bound} indicates that the speed of convergence of the CG method depends on the condition number of the coefficient matrix $A$, i.e., the smaller the condition number is, the faster the CG method converges. It is therefore natural to consider the following equivalent linear systems:}{Theorem\ \ref{thm:cg-condition-number-bound} 表明, CG 方法的收敛速度依赖于系数矩阵 $A$ 的条件数, 即条件数越小, CG 方法收敛越快. 因此很自然地考虑下列等价线性方程组:}
\[A\bm x=\bm b\Leftrightarrow\widetilde A\widetilde{\bm x}=\widetilde{\bm b},\]
\BT{so that the condition number of $\widetilde A$ is much smaller than that of $A$, and then apply the CG method not to $A\bm x=\bm b$ but to the transformed linear systems $\widetilde A\widetilde{\bm x}=\widetilde{\bm b}$. The resulting algorithm is called the preconditioned (PCG) method. The technique to construct $\widetilde A$ is called preconditioning, which will be discussed in Sect.\ \ref{sec:preconditioning-techniques}.}{使得 $\widetilde A$ 的条件数远小于 $A$ 的条件数, 然后不对 $A\bm x=\bm b$, 而对变换后的线性方程组 $\widetilde A\widetilde{\bm x}=\widetilde{\bm b}$ 应用 CG 方法. 所得算法称为预条件 (PCG) 方法. 构造 $\widetilde A$ 的技术称为预条件技术, 将在 Section\ \ref{sec:preconditioning-techniques} 中讨论.}
\BT{We now describe the algorithm of the preconditioned CG method. Let $A\approx K=LL^{\mathrm H}$. Then $\widetilde A=L^{-1}AL^{-\mathrm H}$ is expected to be close to the identity matrix whose condition number is one. We now consider applying the CG method to the following transformed linear system}{现在介绍预条件 CG 方法的算法. 令 $A\approx K=LL^{\mathrm H}$. 则预期 $\widetilde A=L^{-1}AL^{-\mathrm H}$ 接近条件数为一的单位矩阵. 现在考虑对下列变换后的线性方程组应用 CG 方法}
\[\widetilde A\widetilde{\bm x}=\widetilde{\bm b},\]
\BT{where $\widetilde A=L^{-1}AL^{-\mathrm H}$, $\widetilde{\bm x}=L^{\mathrm H}\bm x$, and $\widetilde{\bm b}=L^{-1}\bm b$. Then, the CG method for $\widetilde A\widetilde{\bm x}=\widetilde{\bm b}$ is expected to converge much faster than the CG method for $A\bm x=\bm b$. From Algorithm\ \ref{alg:cg}, the iterates of the CG method for $\widetilde A\widetilde{\bm x}=\widetilde{\bm b}$ are written as}{其中 $\widetilde A=L^{-1}AL^{-\mathrm H}$, $\widetilde{\bm x}=L^{\mathrm H}\bm x$, $\widetilde{\bm b}=L^{-1}\bm b$. 则预期用于 $\widetilde A\widetilde{\bm x}=\widetilde{\bm b}$ 的 CG 方法, 比用于 $A\bm x=\bm b$ 的 CG 方法收敛快得多. 由 Algorithm\ \ref{alg:cg}, 用于 $\widetilde A\widetilde{\bm x}=\widetilde{\bm b}$ 的 CG 方法的迭代式为}
% block: pcg-recurrence-transformation page:94 type:ordinary-paragraph
\[\begin{aligned}
\widetilde{\bm p}_n&=\widetilde{\bm r}_n+\beta_{n-1}\widetilde{\bm p}_{n-1},\\
\alpha_n&=\frac{(\widetilde{\bm r}_n,\widetilde{\bm r}_n)}{(\widetilde{\bm p}_n,\widetilde A\widetilde{\bm p}_n)},\\
\widetilde{\bm x}_{n+1}&=\widetilde{\bm x}_n+\alpha_n\widetilde{\bm p}_n,\\
\widetilde{\bm r}_{n+1}&=\widetilde{\bm r}_n-\alpha_n\widetilde A\widetilde{\bm p}_n,\\
\beta_n&=\frac{(\widetilde{\bm r}_{n+1},\widetilde{\bm r}_{n+1})}{(\widetilde{\bm r}_n,\widetilde{\bm r}_n)}.
\end{aligned}\]
\BT{Using the following rewrites:}{作下列改写:}
\[\widetilde{\bm x}\Rightarrow L^{\mathrm H}\bm x,\quad\widetilde{\bm p}\Rightarrow L^{\mathrm H}\bm p,\quad\widetilde{\bm r}\Rightarrow L^{-1}\bm r,\]
\BT{we have}{得到}
\[\begin{aligned}
\bm p_n&=L^{-\mathrm H}L^{-1}\bm r_n+\beta_{n-1}L^{-\mathrm H}L^{\mathrm H}\bm p_{n-1}=L^{-\mathrm H}L^{-1}\bm r_n+\beta_{n-1}\bm p_{n-1},\\
\alpha_n&=\frac{(L^{-1}\bm r_n,L^{-1}\bm r_n)}{(L^{\mathrm H}\bm p_n,L^{-1}AL^{-\mathrm H}L^{\mathrm H}\bm p_n)}=\frac{(L^{-\mathrm H}L^{-1}\bm r_n,\bm r_n)}{(\bm p_n,A\bm p_n)},\\
\bm x_{n+1}&=L^{-\mathrm H}L^{\mathrm H}\bm x_n+\alpha_nL^{-\mathrm H}L^{\mathrm H}\bm p_n=\bm x_n+\alpha_n\bm p_n,\\
\bm r_{n+1}&=LL^{-1}\bm r_n-\alpha_nLL^{-1}AL^{-\mathrm H}L^{\mathrm H}\bm p_n=\bm r_n-\alpha_nA\bm p_n,\\
\beta_n&=\frac{(L^{-1}\bm r_{n+1},L^{-1}\bm r_{n+1})}{(L^{-\mathrm H}L^{-1}\bm r_n,\bm r_n)}=\frac{(L^{-\mathrm H}L^{-1}\bm r_{n+1},\bm r_{n+1})}{(L^{-\mathrm H}L^{-1}\bm r_n,\bm r_n)}.
\end{aligned}\]
\BT{Since $K^{-1}=L^{-\mathrm H}L^{-1}$, we have the PCG method that is described in Algorithm\ \ref{alg:pcg}.}{由于 $K^{-1}=L^{-\mathrm H}L^{-1}$, 得到 Algorithm\ \ref{alg:pcg} 所述的 PCG 方法.}
% block: pcg-algorithm page:94 type:ordinary-paragraph
\begin{algorithm}[H]\caption{\BT{The preconditioned CG method}{预条件 CG 方法}}\label{alg:pcg}
\begin{algorithmic}[1]
\Statex\textbf{Input:} $\bm x_0\in\mathbb C^N,\ \beta_{-1}=0,\ \bm p_{-1}=\bm0,\ \bm r_0=\bm b-A\bm x_0$
\Statex\textbf{Output:} $\bm x_n$
\For{$n=0,1,2,\ldots$, \BT{until convergence}{直至收敛}}
\State $\bm p_n=K^{-1}\bm r_n+\beta_{n-1}\bm p_{n-1}$
\State $\alpha_n=\frac{(K^{-1}\bm r_n,\bm r_n)}{(\bm p_n,A\bm p_n)}$
\State $\bm x_{n+1}=\bm x_n+\alpha_n\bm p_n$
\State $\bm r_{n+1}=\bm r_n-\alpha_nA\bm p_n$
\State $\beta_n=\frac{(K^{-1}\bm r_{n+1},\bm r_{n+1})}{(K^{-1}\bm r_n,\bm r_n)}$
\EndFor\end{algorithmic}\end{algorithm}
% block: cg-notes pages:94,95 type:ordinary-paragraph
\BT{Some notes on the CG method are listed next.}{下面列出关于 CG 方法的几点说明.}
\begin{enumerate}
\item \BT{From Theorem\ \ref{thm:cg-energy-optimality}, in exact precision arithmetic, the CG method finds the exact solution within $N$ iteration steps.}{由 Theorem\ \ref{thm:cg-energy-optimality}, 在精确算术下, CG 方法在 $N$ 次迭代内找到精确解.}
\item \BT{We can still use the CG method even if the coefficient matrix is not positive definite. In this case, the optimality in Theorem\ \ref{thm:cg-energy-optimality} does not hold anymore, and the CG method may suffer from breakdown, i.e., zero division, because $(\bm p_n,A\bm p_n)>0$ for any $\bm p_n\ne\bm0$ does not hold anymore, see line 3 of Algorithm\ \ref{alg:cg}. In practice, breakdown (zero division) does not occur, but near breakdown may occur, i.e., $(\bm p_n,A\bm p_n)\approx0$, leading to numerical instability.}{即使系数矩阵不正定, 仍可使用 CG 方法. 此时, Theorem\ \ref{thm:cg-energy-optimality} 中的最优性不再成立, CG 方法可能发生中断, 即除以零, 因为对任意 $\bm p_n\ne\bm0$ 都有 $(\bm p_n,A\bm p_n)>0$ 这一性质不再成立, 见 Algorithm\ \ref{alg:cg} 的第 3 行. 实际计算中不会发生中断 (除以零), 但可能发生近中断, 即 $(\bm p_n,A\bm p_n)\approx0$, 导致数值不稳定.}
\item \BT{In exact precision arithmetic, the CG method for Hermitian indefinite matrices also finds the exact solution within $N$ iteration steps if breakdown does not occur.}{在精确算术下, 若不发生中断, 用于 Hermitian 不定矩阵的 CG 方法也在 $N$ 次迭代内找到精确解.}
\end{enumerate}
\BT{In the next two subsections, the CR method and the MINRES method for Hermitian linear systems are described.}{接下来两小节介绍 Hermitian 线性方程组的 CR 方法和 MINRES 方法.}
% block: conjugate-residual-heading page:95 type:subsection
\subsubsection{\BT{The Conjugate Residual (CR) Method}{共轭残差 (CR) 方法}}\label{subsec:conjugate-residual}
\BT{The Conjugate Residual (CR) method \BookCite{177} was proposed by in 1955 Stiefel, who is one of the authors of the CG method. The feature of the CR method is that it has optimality in terms of the residual norm, which will be described later.}{共轭残差 (CR) 方法 \BookCite{177} 由 CG 方法的提出者之一 Stiefel 于 1955 年提出. CR 方法的特点是关于残差范数具有最优性, 稍后将予以说明.}
\BT{It is possible to derive the CR method like the derivation of the CG method, but here we give a concise derivation of the CR method using the algorithm of the CG method. A drawback of the derivation here is that $A$ is assumed to be Hermitian positive definite, whereas the CR method can be used for Hermitian indefinite linear systems. On the other hand, we will see the derivation is much simpler than that of the CG method. A derivation of the CR method similar to that of the CG method can be found in Sect.\ \ref{subsec:cocr-method} by restricting the derivation of the COCR method for complex symmetric linear systems to that for real symmetric ones.}{可以像推导 CG 方法一样推导 CR 方法, 但这里利用 CG 方法的算法给出 CR 方法的简洁推导. 这一推导的缺点是, 假设 $A$ 为 Hermitian 正定矩阵, 而 CR 方法可用于 Hermitian 不定线性方程组. 另一方面, 我们将看到这一推导比 CG 方法的推导简单得多. 与 CG 方法推导相似的 CR 方法推导见 Section\ \ref{subsec:cocr-method}: 将复对称线性方程组的 COCR 方法的推导限制到实对称情形即可得到.}
% block: cr-square-root-transformation page:95 type:ordinary-paragraph
\BT{We now derive the CR method from the CG method. If $A$ is Hermitian positive definite, then there exists a square root of $A$: the Cholesky factorization of $A$, see Sect.\ \ref{subsec:cholesky-decomposition}, is obtained as $LL^{\mathrm H}$, and $L$ can be decomposed as $U\Sigma V^{\mathrm H}$ by using singular value decomposition.}{现在从 CG 方法推导 CR 方法. 若 $A$ 为 Hermitian 正定矩阵, 则 $A$ 存在平方根: $A$ 的 Cholesky 分解 (见 Section\ \ref{subsec:cholesky-decomposition}) 为 $LL^{\mathrm H}$, 利用奇异值分解又可将 $L$ 分解为 $U\Sigma V^{\mathrm H}$.}\footnote{\BT{$U$ and $V$ are unitary matrices, and $\Sigma$ is a diagonal matrix whose diagonal elements are nonnegative.}{$U$ 和 $V$ 为酉矩阵, $\Sigma$ 为对角元素非负的对角矩阵.}}
\BT{Then, we have}{于是有}
\[A=LL^{\mathrm H}=(U\Sigma V^{\mathrm H})(V\Sigma U^{\mathrm H})=U\Sigma^2U^{\mathrm H}=(U\Sigma U^{\mathrm H})(U\Sigma U^{\mathrm H})=A^{\frac12}A^{\frac12}.\]
\BT{Thus, the square root of $A$ is $A^{\frac12}=U\Sigma U^{\mathrm H}$. Here, we show that if $A$ is Hermitian positive definite, then the algorithm of the CR method is obtained by applying the CG method to the following linear systems:}{因此 $A$ 的平方根为 $A^{\frac12}=U\Sigma U^{\mathrm H}$. 这里说明, 若 $A$ 为 Hermitian 正定矩阵, 则对下列线性方程组应用 CG 方法即可得到 CR 方法的算法:}
\begin{equation}\label{eq:cr-square-root-system}A\widetilde{\bm x}=A^{\frac12}\bm b,\quad\widetilde{\bm x}=A^{\frac12}\bm x.\end{equation}
\BT{It follows from the algorithm of the CG method that}{由 CG 方法的算法可得}
% block: cr-transformed-recurrences page:96 type:ordinary-paragraph
\[\begin{aligned}
\widetilde{\bm p}_n&=\widetilde{\bm r}_n+\beta_{n-1}\widetilde{\bm p}_{n-1},\\
\alpha_n&=\frac{(\widetilde{\bm r}_n,\widetilde{\bm r}_n)}{(\widetilde{\bm p}_n,A\widetilde{\bm p}_n)},\\
\widetilde{\bm x}_{n+1}&=\widetilde{\bm x}_n+\alpha_n\widetilde{\bm p}_n,\\
\widetilde{\bm r}_{n+1}&=\widetilde{\bm r}_n-\alpha_nA\widetilde{\bm p}_n,\\
\beta_n&=\frac{(\widetilde{\bm r}_{n+1},\widetilde{\bm r}_{n+1})}{(\widetilde{\bm r}_n,\widetilde{\bm r}_n)}.
\end{aligned}\]
\BT{The residual vector of the original system $A\bm x=\bm b$ is $\bm r_n=\bm b-A\bm x_n$, and then from \eqref{eq:cr-square-root-system} we have $\widetilde{\bm r}_n=A^{\frac12}\bm b-A\widetilde{\bm x}_n=A^{\frac12}(\bm b-A\bm x_n)=A^{\frac12}\bm r_n$. Substituting $\widetilde{\bm r}_n=A^{\frac12}\bm r_n$, $\widetilde{\bm p}_n=A^{\frac12}\bm p_n$, and $\widetilde{\bm x}_n=A^{\frac12}\bm x_n$ into the above recurrences, we have the algorithm of the CR method (Algorithm\ \ref{alg:cr}).}{原方程组 $A\bm x=\bm b$ 的残差向量为 $\bm r_n=\bm b-A\bm x_n$, 由 \eqref{eq:cr-square-root-system} 有 $\widetilde{\bm r}_n=A^{\frac12}\bm b-A\widetilde{\bm x}_n=A^{\frac12}(\bm b-A\bm x_n)=A^{\frac12}\bm r_n$. 将 $\widetilde{\bm r}_n=A^{\frac12}\bm r_n$, $\widetilde{\bm p}_n=A^{\frac12}\bm p_n$, $\widetilde{\bm x}_n=A^{\frac12}\bm x_n$ 代入上述递推式, 得到 CR 方法的算法 (Algorithm\ \ref{alg:cr}).}
% block: cr-algorithm page:96 type:ordinary-paragraph
\begin{algorithm}[H]\caption{\BT{The CR method}{CR 方法}}\label{alg:cr}\begin{algorithmic}[1]
\Statex\textbf{Input:} $\bm x_0\in\mathbb C^N,\ \beta_{-1}=0,\ \bm p_{-1}=\bm0,\ \bm r_0=\bm b-A\bm x_0$
\Statex\textbf{Output:} $\bm x_n$
\For{$n=0,1,2,\ldots$, \BT{until convergence}{直至收敛}}
\State $\bm p_n=\bm r_n+\beta_{n-1}\bm p_{n-1}$
\State $A\bm p_n=A\bm r_n+\beta_{n-1}A\bm p_{n-1}$
\State $\alpha_n=\frac{(\bm r_n,A\bm r_n)}{(A\bm p_n,A\bm p_n)}$
\State $\bm x_{n+1}=\bm x_n+\alpha_n\bm p_n$
\State $\bm r_{n+1}=\bm r_n-\alpha_nA\bm p_n$
\State $\beta_n=\frac{(\bm r_{n+1},A\bm r_{n+1})}{(\bm r_n,A\bm r_n)}$
\EndFor\end{algorithmic}\end{algorithm}
% block: cr-orthogonality-and-optimality pages:96,97 type:ordinary-paragraph
\BT{From line 7 of Algorithm\ \ref{alg:cr}, the CR method may suffer from breakdown if Hermitian matrix $A$ is indefinite, i.e., $(\bm r_n,A\bm r_n)=0$. By induction, it can be shown that the CR method generates iterates $\bm r_i,\bm p_i$ that satisfy}{由 Algorithm\ \ref{alg:cr} 的第 7 行, 若 Hermitian 矩阵 $A$ 不定, CR 方法可能发生中断, 即 $(\bm r_n,A\bm r_n)=0$. 利用归纳法可证明, CR 方法生成的迭代向量 $\bm r_i,\bm p_i$ 满足}
\begin{equation}\label{eq:cr-residual-orthogonality}(\bm r_i,A\bm r_j)=0\quad\text{for }i\ne j,\end{equation}
\begin{equation}\label{eq:cr-direction-orthogonality}(A\bm p_i,A\bm p_j)=0\quad\text{for }i\ne j.\end{equation}
\BT{Comparing the two properties (\eqref{eq:cr-residual-orthogonality} and \eqref{eq:cr-direction-orthogonality}) and those of the CG method, the CR method finds approximate solutions such that}{将两个性质 (\eqref{eq:cr-residual-orthogonality} 和 \eqref{eq:cr-direction-orthogonality}) 与 CG 方法的性质比较, CR 方法寻找满足下列条件的近似解}
\[\bm x_n=\bm x_0+\bm z_n,\quad\bm z_n\in\mathcal K_n(A,\bm r_0),\]
\BT{where $\bm z_n$ is determined by $\bm r_n\perp A\mathcal K_n(A,\bm r_0)(:=\mathcal K_n(A,A\bm r_0))$, or equivalently}{其中 $\bm z_n$ 由 $\bm r_n\perp A\mathcal K_n(A,\bm r_0)(:=\mathcal K_n(A,A\bm r_0))$ 确定, 或等价地,}
\[\bm x_n=\bm x_0+V_n\bm y_n,\quad\bm y_n\in\mathbb C^n,\]
\BT{where $\bm y_n$ is determined by $(AV_n)^{\mathrm H}\bm r_n=\bm0$. Following the proof of Theorem\ \ref{thm:cg-energy-optimality}, we have}{其中 $\bm y_n$ 由 $(AV_n)^{\mathrm H}\bm r_n=\bm0$ 确定. 沿用 Theorem\ \ref{thm:cg-energy-optimality} 的证明, 得}
\[\min_{\bm x_n\in\bm x_0+\mathcal K_n(A,\bm r_0)}\|\bm e_n\|_{A^2}.\]
\BT{Since $\|\bm e_n\|_{A^2}^2=\bm e_n^{\mathrm H}A^2\bm e_n=(A\bm e_n)^{\mathrm H}A\bm e_n=\bm r_n^{\mathrm H}\bm r_n=\|\bm r_n\|^2$, the CR method finds approximate solutions such that}{由于 $\|\bm e_n\|_{A^2}^2=\bm e_n^{\mathrm H}A^2\bm e_n=(A\bm e_n)^{\mathrm H}A\bm e_n=\bm r_n^{\mathrm H}\bm r_n=\|\bm r_n\|^2$, CR 方法寻找满足下列条件的近似解}
\[\min_{\bm x_n\in\bm x_0+\mathcal K_n(A,\bm r_0)}\|\bm r_n\|.\]
\BT{Therefore, the CR method produces the optimal approximate solution in terms of the residual norm. Similar to the derivation of the preconditioned CG method, the preconditioned CR method can be derived from applying the CR method to the transformed linear systems $\widetilde A\widetilde{\bm x}=\widetilde{\bm b}$. The resulting algorithm is described in Algorithm\ \ref{alg:pcr}.}{因此, CR 方法产生关于残差范数的最优近似解. 与预条件 CG 方法的推导类似, 对变换后的线性方程组 $\widetilde A\widetilde{\bm x}=\widetilde{\bm b}$ 应用 CR 方法, 可推导出预条件 CR 方法. 所得算法见 Algorithm\ \ref{alg:pcr}.}
% block: pcr-algorithm page:97 type:ordinary-paragraph
\begin{algorithm}[H]\caption{\BT{The preconditioned CR method}{预条件 CR 方法}}\label{alg:pcr}\begin{algorithmic}[1]
\Statex\textbf{Input:} $\bm x_0\in\mathbb C^N,\ \beta_{-1}=0,\ \bm p_{-1}=\bm0,\ \bm r_0=\bm b-A\bm x_0$
\Statex\textbf{Output:} $\bm x_n$
\For{$n=0,1,2,\ldots$, \BT{until convergence}{直至收敛}}
\State $\bm p_n=K^{-1}\bm r_n+\beta_{n-1}\bm p_{n-1}$
\State $(A\bm p_n=AK^{-1}\bm r_n+\beta_{n-1}A\bm p_{n-1})$
\State $\alpha_n=\frac{(K^{-1}\bm r_n,AK^{-1}\bm r_n)}{(A\bm p_n,K^{-1}A\bm p_n)}$
\State $\bm x_{n+1}=\bm x_n+\alpha_n\bm p_n$
\State $\bm r_{n+1}=\bm r_n-\alpha_nA\bm p_n$
\State $(K^{-1}\bm r_{n+1}=K^{-1}\bm r_n-\alpha_nK^{-1}A\bm p_n)$
\State $\beta_n=\frac{(K^{-1}\bm r_{n+1},AK^{-1}\bm r_{n+1})}{(K^{-1}\bm r_n,AK^{-1}\bm r_n)}$
\EndFor\end{algorithmic}\end{algorithm}
% block: pcr-operation-count page:97 type:ordinary-paragraph
\BT{Note that lines 3 and 7 in Algorithm\ \ref{alg:pcr} are added for reducing the number of matrix--vector multiplications of the form $A\bm v$ and solving linear systems of the form $K\bm v=\bm w$. In fact, at each iteration, it seems that we need to compute}{注意, Algorithm\ \ref{alg:pcr} 中加入第 3 和第 7 行, 是为了减少形如 $A\bm v$ 的矩阵向量乘法及形如 $K\bm v=\bm w$ 的线性方程组求解次数. 实际上, 每次迭代似乎都需要计算}
\[A\bm p_n,\quad K^{-1}\bm r_{n+1},\quad A(K^{-1}\bm r_n),\quad K^{-1}(A\bm p_n).\]
\BT{But from lines 3 and 7, we do not need to compute $A\bm p_n$ and $K^{-1}\bm r_{n+1}$, except $K^{-1}\bm r_0$ at the first iteration step. Thus, essentially we only need to compute the following form:}{但由第 3 和第 7 行, 除第一次迭代需要计算 $K^{-1}\bm r_0$ 外, 不必计算 $A\bm p_n$ 和 $K^{-1}\bm r_{n+1}$. 因此, 每次迭代实质上只需计算下列形式:}
\[A\bm v,\quad K^{-1}\bm w\]
\BT{at each iteration step. Here, $\bm v:=K^{-1}\bm r_{n+1}$ is obtained by line 7 and $\bm w:=A\bm p_n$ is obtained by line 3.}{这里 $\bm v:=K^{-1}\bm r_{n+1}$ 由第 7 行得到, $\bm w:=A\bm p_n$ 由第 3 行得到.}
% block: minres-heading page:98 type:subsection
\subsubsection{\BT{The Minimal Residual (MINRES) Method}{最小残差 (MINRES) 方法}}\label{subsec:minres}
\BT{The Minimal Residual (MINRES) method \BookCite{143} was proposed by Paige and Saunders in 1975. The MINRES method generates $\bm x_n$ that minimizes $\|\bm b-A\bm x_n\|$ over the affine space $\bm x_0+\mathcal K_n(A,\bm r_0)$. Therefore, in exact precision arithmetic, the MINRES method and the CR method produce the same approximate solutions. The minimization can be achieved by using the Lanczos process.}{最小残差 (MINRES) 方法 \BookCite{143} 由 Paige 和 Saunders 于 1975 年提出. MINRES 方法生成在仿射空间 $\bm x_0+\mathcal K_n(A,\bm r_0)$ 上使 $\|\bm b-A\bm x_n\|$ 达到极小值的 $\bm x_n$. 因此在精确算术下, MINRES 方法与 CR 方法产生相同的近似解. 这一极小化可利用 Lanczos 过程实现.}
\BT{Here, we give a derivation process of the MINRES method. Let $V_n$ be the orthonormal basis of $\mathcal K_n(A,\bm r_0)$. The MINRES method finds $\bm x_n$ over the affine space $\bm x_0+\mathcal K_n(A,\bm r_0)$, i.e.,}{这里给出 MINRES 方法的推导过程. 设 $V_n$ 为 $\mathcal K_n(A,\bm r_0)$ 的标准正交基. MINRES 方法在仿射空间 $\bm x_0+\mathcal K_n(A,\bm r_0)$ 上寻找 $\bm x_n$, 即}
\begin{equation}\label{eq:minres-affine-solution}\bm x_n=\bm x_0+V_n\bm y_n,\quad\bm y_n\in\mathbb C^n.\end{equation}
\BT{The corresponding residual vector is given by}{相应的残差向量为}
\[\bm r_n=\bm r_0-AV_n\bm y_n.\]
\BT{From the matrix form of the Lanczos process in \eqref{eq:lanczos-matrix-relation}, it follows that}{由 \eqref{eq:lanczos-matrix-relation} 中 Lanczos 过程的矩阵形式, 得}
\[\bm r_n=\bm r_0-V_{n+1}T_{n+1,n}\bm y_n=V_{n+1}(\beta\bm e_1-T_{n+1,n}\bm y_n),\]
\BT{where $\beta:=\|\bm r_0\|$. The 2-norm of the residual vector is written as}{其中 $\beta:=\|\bm r_0\|$. 残差向量的 2-范数写成}
\[\begin{aligned}
\|\bm r_n\|&=\|V_{n+1}(\beta\bm e_1-T_{n+1,n}\bm y_n)\|\\
&=\sqrt{[V_{n+1}(\beta\bm e_1-T_{n+1,n}\bm y_n)]^{\mathrm H}V_{n+1}(\beta\bm e_1-T_{n+1,n}\bm y_n)}\\
&=\sqrt{(\beta\bm e_1-T_{n+1,n}\bm y_n)^{\mathrm H}V_{n+1}^{\mathrm H}V_{n+1}(\beta\bm e_1-T_{n+1,n}\bm y_n)}\\
&=\sqrt{(\beta\bm e_1-T_{n+1,n}\bm y_n)^{\mathrm H}(\beta\bm e_1-T_{n+1,n}\bm y_n)}\\
&=\|\beta\bm e_1-T_{n+1,n}\bm y_n\|.
\end{aligned}\]
\BT{The MINRES method finds approximate solutions such that the 2-norm of the residual vector is minimized, i.e., $\bm y_n$ is chosen such that}{MINRES 方法寻找使残差向量的 2-范数达到极小值的近似解, 即选择 $\bm y_n$ 使得}
\begin{equation}\label{eq:minres-projected-least-squares}\bm y_n:=\operatorname*{arg\,min}_{\bm y\in\mathbb C^n}\|\beta\bm e_1-T_{n+1,n}\bm y\|.\end{equation}
\BT{By solving the least-squares problem in \eqref{eq:minres-projected-least-squares} using a Givens rotation that is one of the unitary matrices, the MINRES method is obtained. In the following, we describe how Givens rotations work for solving \eqref{eq:minres-projected-least-squares}. For simplicity, we consider the case $n=4$ for solving it. $\beta\bm e_1-T_{5,4}\bm y$ is written as}{利用作为酉矩阵的一种 Givens 旋转求解 \eqref{eq:minres-projected-least-squares} 中的最小二乘问题, 就得到 MINRES 方法. 下面说明 Givens 旋转如何用于求解 \eqref{eq:minres-projected-least-squares}. 为简单起见, 考虑 $n=4$ 的情形. $\beta\bm e_1-T_{5,4}\bm y$ 写成}
% block: minres-first-givens page:99 type:ordinary-paragraph
\[\beta\bm e_1-T_{5,4}\bm y=\begin{bmatrix}\beta\\0\\0\\0\\0\end{bmatrix}
-\begin{bmatrix}t_{11}&t_{12}&&\\t_{21}&t_{22}&t_{23}&\\&t_{32}&t_{33}&t_{34}\\&&t_{43}&t_{44}\\&&&t_{54}\end{bmatrix}\begin{bmatrix}y_1\\y_2\\y_3\\y_4\end{bmatrix},\]
\BT{where $t_{ij}=t_{ji}$, $1\le i,j\le4$. Let $G_1$ be a matrix of the Givens rotation, i.e.,}{其中 $t_{ij}=t_{ji}$, $1\le i,j\le4$. 设 $G_1$ 为 Givens 旋转矩阵, 即}
\[G_1=\begin{bmatrix}c_1&s_1&&&\\-\bar s_1&c_1&&&\\&&1&&\\&&&1&\\&&&&1\end{bmatrix}
\quad\text{with }c_1=\frac{|t_{11}|}{\sqrt{|t_{11}|^2+|t_{21}|^2}},\quad\bar s_1=\frac{t_{21}}{t_{11}}c_1,\]
\BT{and $s_1$ is obtained by the conjugate of $\bar s_1$. There is one exception: set $s_1=\bar s_1=1$ if $t_{11}=0$. Then, $G_1$ is a unitary matrix, and thus $\|\beta\bm e_1-T_{n+1,n}\bm y\|=\|G_1(\beta\bm e_1-T_{n+1,n}\bm y)\|$. This leads to}{其中 $s_1$ 取为 $\bar s_1$ 的共轭. 有一种例外: 若 $t_{11}=0$, 则令 $s_1=\bar s_1=1$. 于是 $G_1$ 为酉矩阵, 因而 $\|\beta\bm e_1-T_{n+1,n}\bm y\|=\|G_1(\beta\bm e_1-T_{n+1,n}\bm y)\|$. 由此得到}
\[G_1(\beta\bm e_1-T_{5,4}\bm y)=\begin{bmatrix}g_1^{(1)}\\g_2^{(1)}\\0\\0\\0\end{bmatrix}
-\begin{bmatrix}t_{11}^{(1)}&t_{12}^{(1)}&t_{13}^{(1)}&\\&t_{22}^{(1)}&t_{23}^{(1)}&\\&t_{32}&t_{33}&t_{34}\\&&t_{43}&t_{44}\\&&&t_{54}\end{bmatrix}\begin{bmatrix}y_1\\y_2\\y_3\\y_4\end{bmatrix},\]
% block: minres-second-givens page:99 type:ordinary-paragraph
\BT{where $g_2^{(1)}=-\bar s_1\beta$. Next, let $G_2$ be a unitary matrix defined by}{其中 $g_2^{(1)}=-\bar s_1\beta$. 接着定义酉矩阵 $G_2$ 为}
\[G_2=\begin{bmatrix}1&&&&\\&c_2&s_2&&\\&-\bar s_2&c_2&&\\&&&1&\\&&&&1\end{bmatrix}
\quad\text{with }c_2=\frac{|t_{22}^{(1)}|}{\sqrt{|t_{22}^{(1)}|^2+|t_{32}|^2}},\quad\bar s_2=\frac{t_{32}}{t_{22}^{(1)}}c_2.\]
\BT{There is one exception: set $s_2=\bar s_2=1$ if $t_{22}^{(1)}=0$. Then, $G_2$ and $G_2G_1$ are unitary matrices, and thus $\|\beta\bm e_1-T_{n+1,n}\bm y\|=\|G_2G_1(\beta\bm e_1-T_{n+1,n}\bm y)\|$. This leads to}{有一种例外: 若 $t_{22}^{(1)}=0$, 则令 $s_2=\bar s_2=1$. 于是 $G_2$ 和 $G_2G_1$ 为酉矩阵, 因而 $\|\beta\bm e_1-T_{n+1,n}\bm y\|=\|G_2G_1(\beta\bm e_1-T_{n+1,n}\bm y)\|$. 由此得到}
\[G_2G_1(\beta\bm e_1-T_{5,4}\bm y)=\begin{bmatrix}g_1^{(1)}\\g_2^{(2)}\\g_3^{(2)}\\0\\0\end{bmatrix}
-\begin{bmatrix}t_{11}^{(1)}&t_{12}^{(1)}&t_{13}^{(1)}&\\&t_{22}^{(2)}&t_{23}^{(2)}&t_{24}^{(2)}\\&&t_{33}^{(2)}&t_{34}^{(2)}\\&&t_{43}&t_{44}\\&&&t_{54}\end{bmatrix}\begin{bmatrix}y_1\\y_2\\y_3\\y_4\end{bmatrix},\]
\BT{where $g_3^{(2)}=(-1)^2\bar s_2\bar s_1\beta$. Similarly, using $G_3$ and $G_4$, we finally obtain}{其中 $g_3^{(2)}=(-1)^2\bar s_2\bar s_1\beta$. 类似地, 利用 $G_3$ 和 $G_4$, 最终得到}
% block: minres-triangular-problem page:100 type:ordinary-paragraph
\[\begin{aligned}
Q_4(\beta\bm e_1-T_{5,4}\bm y)&=\begin{bmatrix}g_1^{(1)}\\g_2^{(2)}\\g_3^{(3)}\\g_4^{(4)}\\g_5^{(5)}\end{bmatrix}
-\begin{bmatrix}t_{11}^{(1)}&t_{12}^{(1)}&t_{13}^{(1)}&\\&t_{22}^{(2)}&t_{23}^{(2)}&t_{24}^{(2)}\\&&t_{33}^{(3)}&t_{34}^{(3)}\\&&&t_{44}^{(4)}\\&&&\end{bmatrix}\begin{bmatrix}y_1\\y_2\\y_3\\y_4\end{bmatrix}\\
&=\begin{bmatrix}\bm g_4\\g_5^{(5)}\end{bmatrix}-\begin{bmatrix}R_4\\\bm0^{\top}\end{bmatrix}\bm y_4,
\end{aligned}\]
\BT{where $Q_4=G_4G_3G_2G_1$ and $g_5^{(5)}=(-1)^4\bar s_4\bar s_3\bar s_2\bar s_1\beta$. Hence, we obtain $\bm y_4$ by solving the following equation:}{其中 $Q_4=G_4G_3G_2G_1$, $g_5^{(5)}=(-1)^4\bar s_4\bar s_3\bar s_2\bar s_1\beta$. 因此, 通过回代求解下列方程, 得到 $\bm y_4$:}
\begin{equation}\label{eq:minres-upper-triangular-system}R_4\bm y_4=\bm g_4\end{equation}
\BT{using backward substitution. It is interesting to see that we can obtain the value of $\min_{\bm y\in\mathbb C^4}\|\beta\bm e_1-T_{5,4}\bm y\|$ without solving the equation $R_4\bm y_4=\bm g_4$, since the value satisfies}{有趣的是, 不必求解方程 $R_4\bm y_4=\bm g_4$, 即可得到 $\min_{\bm y\in\mathbb C^4}\|\beta\bm e_1-T_{5,4}\bm y\|$ 的值, 因为该值满足}
\[\begin{aligned}
\|\bm r_4\|&=\min_{\bm y\in\mathbb C^4}\|\beta\bm e_1-T_{5,4}\bm y\|
=\min_{\bm y\in\mathbb C^4}\|Q_4(\beta\bm e_1-T_{5,4}\bm y)\|\\
&=\left\|\begin{bmatrix}\bm g_4\\g_5^{(5)}\end{bmatrix}-\begin{bmatrix}R_4\\\bm0^{\top}\end{bmatrix}\bm y_4\right\|
=\left\|\begin{bmatrix}\bm0\\g_5^{(5)}\end{bmatrix}\right\|\\
&=|g_5^{(5)}|.
\end{aligned}\]
% block: minres-solution-recurrences pages:100,101 type:ordinary-paragraph
\BT{Thus, $|g_5^{(5)}|$ corresponds to the residual norm of $\bm r_4$, and this can be used as a stopping criterion. To obtain the approximate solution, it follows from \eqref{eq:minres-affine-solution} and \eqref{eq:minres-upper-triangular-system} that}{因此, $|g_5^{(5)}|$ 对应于 $\bm r_4$ 的残差范数, 可用作停止准则. 为得到近似解, 由 \eqref{eq:minres-affine-solution} 和 \eqref{eq:minres-upper-triangular-system} 有}
\begin{equation}\label{eq:minres-solution-triangular-inverse}\bm x_4=\bm x_0+V_4R_4^{-1}\bm g_4.\end{equation}
\BT{Here, we introduce a matrix $P_4:=V_4R_4^{-1}$ with columns $[\bm p_1,\bm p_2,\bm p_3,\bm p_4]$. Then, from $[\bm p_1,\bm p_2,\bm p_3,\bm p_4]R_4=[\bm v_1,\bm v_2,\bm v_3,\bm v_4]$, we have the following recurrences:}{这里引入矩阵 $P_4:=V_4R_4^{-1}$, 其列为 $[\bm p_1,\bm p_2,\bm p_3,\bm p_4]$. 于是由 $[\bm p_1,\bm p_2,\bm p_3,\bm p_4]R_4=[\bm v_1,\bm v_2,\bm v_3,\bm v_4]$, 有下列递推关系:}
\begin{equation}\label{eq:minres-direction-one}\bm p_1=\bm v_1/t_{11}^{(1)},\end{equation}
\begin{equation}\label{eq:minres-direction-two}\bm p_2=(\bm v_2-t_{12}^{(1)}\bm p_1)/t_{22}^{(2)},\end{equation}
\begin{equation}\label{eq:minres-direction-three}\bm p_3=(\bm v_3-t_{13}^{(1)}\bm p_1-t_{23}^{(2)}\bm p_2)/t_{33}^{(3)},\end{equation}
\begin{equation}\label{eq:minres-direction-four}\bm p_4=(\bm v_4-t_{24}^{(2)}\bm p_2-t_{34}^{(3)}\bm p_3)/t_{44}^{(4)}.\end{equation}
\BT{From \eqref{eq:minres-solution-triangular-inverse}, $\bm x_4=\bm x_0+P_4\bm g_4=\bm x_0+\sum_{i=1}^4g_i\bm p_i$ and thus}{由 \eqref{eq:minres-solution-triangular-inverse}, $\bm x_4=\bm x_0+P_4\bm g_4=\bm x_0+\sum_{i=1}^4g_i\bm p_i$, 因此}
\begin{equation}\label{eq:minres-solution-updates}\bm x_i=\bm x_{i-1}+g_i\bm p_i,\quad i=1,\ldots,4,\end{equation}
\BT{where $g_i$ is the $i$th element of $\bm g_4$. The approximate solutions $\bm x_1,\ldots,\bm x_4$ of the MINRES method are obtained by \eqref{eq:minres-solution-updates} with the recurrences \eqref{eq:minres-direction-one}--\eqref{eq:minres-direction-four}. The derivation for the case $n=4$ can be easily generalized to the $n$th iteration step, leading to Algorithm\ \ref{alg:minres}. From Algorithm\ \ref{alg:minres}, the MINRES method never suffers from breakdown since $\beta_n\ne0$ unless $\widetilde{\bm v}_{n+1}=\bm0$.}{其中 $g_i$ 是 $\bm g_4$ 的第 $i$ 个元素. MINRES 方法的近似解 $\bm x_1,\ldots,\bm x_4$ 由 \eqref{eq:minres-solution-updates} 及递推式 \eqref{eq:minres-direction-one}--\eqref{eq:minres-direction-four} 得到. $n=4$ 时的推导容易推广到第 $n$ 次迭代, 得到 Algorithm\ \ref{alg:minres}. 由 Algorithm\ \ref{alg:minres}, MINRES 方法从不发生中断, 因为除非 $\widetilde{\bm v}_{n+1}=\bm0$, 否则 $\beta_n\ne0$.}
% block: minres-algorithm page:101 type:ordinary-paragraph
\begin{algorithm}[H]\caption{\BT{The MINRES method}{MINRES 方法}}\label{alg:minres}\begin{algorithmic}[1]
\Statex\textbf{Input:} $\bm x_0\in\mathbb C^N,\ \beta_0=0,\ \bm v_0=\bm0,\ \bm r_0=\bm b-A\bm x_0$
\Statex\textbf{Output:} $\bm x_n$
\State $\bm g=(\|\bm r_0\|,0,\ldots,0)^{\top},\ \bm v_1=\bm r_0/\|\bm r_0\|$
\For{$n=1,2,\ldots$}
\State (Lanczos \BT{process}{过程})
\State $\alpha_n=(\bm v_n,A\bm v_n)$
\State $\widetilde{\bm v}_{n+1}=A\bm v_n-\alpha_n\bm v_n-\beta_{n-1}\bm v_{n-1}$
\State $\beta_n=(\widetilde{\bm v}_{n+1},\widetilde{\bm v}_{n+1})^{1/2}$
\State $\bm v_{n+1}=\widetilde{\bm v}_{n+1}/\beta_n$
\State $t_{n-1,n}=\beta_{n-1},\ t_{n,n}=\alpha_n$
\State $t_{n+1,n}=\beta_n$
\State (Givens \BT{rotations}{旋转})
\For{$i=\max\{1,n-2\},\ldots,n-1$}
\State $\begin{bmatrix}t_{i,n}\\t_{i+1,n}\end{bmatrix}=\begin{bmatrix}c_i&s_i\\-\bar s_i&c_i\end{bmatrix}\begin{bmatrix}t_{i,n}\\t_{i+1,n}\end{bmatrix}$
\EndFor
\State $c_n=\frac{|t_{n,n}|}{\sqrt{|t_{n,n}|^2+|t_{n+1,n}|^2}}$
\State $\bar s_n=\frac{t_{n+1,n}}{t_{n,n}}c_n$
\State $t_{n,n}=c_nt_{n,n}+s_nt_{n+1,n}$
\State $t_{n+1,n}=0$
\State $\begin{bmatrix}g_n\\g_{n+1}\end{bmatrix}=\begin{bmatrix}c_n&s_n\\-\bar s_n&c_n\end{bmatrix}\begin{bmatrix}g_n\\0\end{bmatrix}$
\State (\BT{Update}{更新} $\bm x_n$)
\State $\bm p_n=(\bm v_n-t_{n-2,n}\bm p_{n-2}$
\State $\qquad-t_{n-1,n}\bm p_{n-1})/t_{n,n}$
\State $\bm x_n=\bm x_{n-1}+g_n\bm p_n$
\State (\BT{Check convergence}{检查收敛})
\State \BT{if $|g_{n+1}|/\|\bm b\|\le\epsilon$, then stop}{若 $|g_{n+1}|/\|\bm b\|\le\epsilon$, 则停止}
\EndFor\end{algorithmic}\end{algorithm}
% block: symmlq-projected-system pages:101,102 type:ordinary-paragraph
\BT{As mentioned above, the MINRES method never suffers from breakdown. Paige and Saunders also proposed the SYMMLQ method in \BookCite{143}, which is a variant of the CG method, to avoid breakdown for Hermitian indefinite linear systems. A brief explanation is given next. Recall that the approximate solution $\bm x_n$ of the CG method is given by}{如上所述, MINRES 方法从不发生中断. Paige 和 Saunders 还在 \BookCite{143} 中提出了 CG 方法的一种变体 SYMMLQ 方法, 以避免 Hermitian 不定线性方程组求解中的中断. 下面作简要说明. 回顾 CG 方法的近似解 $\bm x_n$ 为}
\[\begin{aligned}
\bm x_n&=\bm x_0+V_n\bm y_n,\\
\bm r_n&=\bm r_0-V_{n+1}T_{n+1,n}\bm y_n=V_{n+1}(\beta\bm e_1-T_{n+1,n}\bm y_n),
\end{aligned}\]
\BT{where $\bm y_n$ is determined by $\bm r_n\perp\mathcal K_n(A,\bm r_0)$ (see Corollary\ \ref{cor:cg-krylov-orthogonality}), i.e., $V^{\mathrm H}\bm r_n=\bm0$. This yields}{其中 $\bm y_n$ 由 $\bm r_n\perp\mathcal K_n(A,\bm r_0)$ 确定 (见 Corollary\ \ref{cor:cg-krylov-orthogonality}), 即 $V^{\mathrm H}\bm r_n=\bm0$. 由此得到}
\begin{equation}\label{eq:symmlq-projected-system}T_n\bm y_n=\beta\bm e_1,\end{equation}
\BT{where $T_n$ is the tridiagonal matrix obtained by removing the last row of $T_{n+1,n}$. The CG method is obtained by the $LDL^{\mathrm H}$ decomposition of $T_n$, i.e.,}{其中 $T_n$ 是删除 $T_{n+1,n}$ 最后一行得到的三对角矩阵. CG 方法通过 $T_n$ 的 $LDL^{\mathrm H}$ 分解得到, 即}
\[T_n=LDL^{\mathrm H},\]
\BT{where $L$ is a unit lower bidiagonal matrix and $D$ is a diagonal matrix. If $A$ is Hermitian positive definite, $T_n$ is also Hermitian positive definite. Thus this decomposition never suffers from breakdown. On the other hand, when $A$ is indefinite, $LDL^{\mathrm H}$ of $T_n$ may not be obtained due to breakdown.}{其中 $L$ 为单位下双对角矩阵, $D$ 为对角矩阵. 若 $A$ 为 Hermitian 正定矩阵, 则 $T_n$ 也是 Hermitian 正定矩阵. 因此这一分解从不发生中断. 另一方面, 当 $A$ 不定时, 可能因中断而无法得到 $T_n$ 的 $LDL^{\mathrm H}$ 分解.}
% block: symmlq-lq-factorization page:102 type:ordinary-paragraph
\BT{In order to circumvent the problem, Paige and Saunders proposed to use the LQ factorization of $T_n$, i.e.,}{为绕过这一问题, Paige 和 Saunders 提议使用 $T_n$ 的 LQ 分解, 即}
\[T_n=L_nQ_n,\quad Q_n^{\mathrm H}Q_n=I,\]
\BT{where $L_n$ is a lower triangular matrix and the decomposition is obtained by Givens rotations. Thus the decomposition is free from breakdown. The resulting algorithm is known as the SYMMLQ method. When the coefficient matrix $A$ is Hermitian positive definite, the SYMMLQ method gives the same result as the CG method in exact precision arithmetic, and for indefinite cases, the SYMMLQ method succeeded in avoiding breakdown when solving \eqref{eq:symmlq-projected-system}. For a brief historical note including the relation between the SYMMLQ method and (less-known) Fridman's method of 1962, see Sect.\ 2.4 in \BookCite{160}.}{其中 $L_n$ 为下三角矩阵, 该分解通过 Givens 旋转得到, 因而不发生中断. 所得算法称为 \linebreak[4]SYMMLQ 方法. 当系数矩阵 $A$ 为 Hermitian 正定矩阵时, 在精确算术下 SYMMLQ 方法给出与 CG 方法相同的结果; 在不定情形下, SYMMLQ 方法在求解 \eqref{eq:symmlq-projected-system} 时成功避免了中断. 关于简要历史说明, 包括 SYMMLQ 方法与 1962 年 (较少为人所知的) Fridman 方法之间的关系, 见 \BookCite{160} 的 Section\ 2.4.}
% block: composite-step-factorization page:102 type:ordinary-paragraph
\BT{Another important remedy is ``composite step'' technique by Bank and Chan \BookCite{16} and similar techniques (hyperbolic pairs) in \BookCite{58}, \BookCite{126}. The key idea of the composite step is decomposing $T_n$ in \eqref{eq:symmlq-projected-system} into}{另一种重要的补救方法是 Bank 和 Chan 的 ``复合步'' 技术 \BookCite{16}, 以及 \BookCite{58}, \BookCite{126} 中的类似技术 (双曲对). 复合步的关键思想是将 \eqref{eq:symmlq-projected-system} 中的 $T_n$ 分解为}
\[T_n=\widetilde L\widetilde D\widetilde L^{\mathrm H},\]
\BT{where $\widetilde L$ is a unit lower block bidiagonal matrix and $\widetilde D$ is a block diagonal matrix whose block is of size $1\times1$ or $2\times2$, i.e.,}{其中 $\widetilde L$ 为单位下分块双对角矩阵, $\widetilde D$ 为块大小为 $1\times1$ 或 $2\times2$ 的分块对角矩阵, 即}
\[\widetilde D=\begin{bmatrix}D_1&&\\&\ddots&\\&&D_m\end{bmatrix},\]
\BT{where $D_i$ is a scalar or a $2\times2$ matrix. This decomposition is also free from breakdown, see, e.g., \BookCite{16}, Theorem 2.2. An algorithmic explanation of the composite step technique (the composite step BiCG method) is described in Sect.\ \ref{subsec:composite-step-bicg}.}{其中 $D_i$ 为标量或 $2\times2$ 矩阵. 这一分解也不发生中断, 例如见 \BookCite{16} 的 Theorem 2.2. 复合步技术的算法说明 (复合步 BiCG 方法) 见 Section\ \ref{subsec:composite-step-bicg}.}
